\documentclass[12pt]{article}
\usepackage[utf8]{inputenc}
\usepackage[T1]{fontenc}
\usepackage{hyperref}
\title{Yilan Creole Language Resources for K-12 Curriculum Development}
\author{Compiled by Claude Sonnet 5.5 (claude-sonnet-5-5) for Andr\'as Kornai}
\date{}
\begin{document}
\maketitle

\section{Introduction}
Yilan Creole (ISO 639-3 ycr; Glottocode yila1234) is a Japanese-based creole spoken in Yilan County, Taiwan. It arose in the 1930s and 1940s from contact between Japanese and Atayal speakers. Sources disagree on speaker numbers: Wikipedia says about 3,000, while TaiwanPlus says fewer than 1,000 and Global Voices says a small number of elders. Global Voices reports that the language lacks official recognition and is not taught in schools.

\section{Primary Educational Resources}
No primers or textbooks were found. No school programme was found.

\section{Linguistic and Reference Materials}
\begin{itemize}
\item Wikipedia, Yilan Creole: \url{https://en.wikipedia.org/wiki/Yilan_Creole}
\item Wikipedia, Yilan Creole Japanese: \url{https://en.wikipedia.org/wiki/Yilan_Creole_Japanese}
\item Peng Qiu, MA thesis (University of Alberta, 2015): \url{https://era.library.ualberta.ca/files/02870z43h/Qiu_Peng_201502_MA.pdf}
\item ``The language endangerment of Yilan Creole in the postcolonial Atayal community of Taiwan'' (ScholarSpace, University of Hawai\textquoteleft i): \url{https://scholarspace.manoa.hawaii.edu/items/84d4b5a9-7814-45e4-ae6b-b28a63f44fdf}
\item The Language Closet, ``When Japanese Met Atayal'' (popular overview): \url{https://thelanguagecloset.com/2020/11/28/when-japanese-met-atayal-yilan-creole-japanese/}
\end{itemize}

\section{Religious and Community Resources}
None were found in the search.

\section{Digital and Audio Resources}
\begin{itemize}
\item Global Voices, ``Taiwan's Yilan Creole: Revitalizing the only Japanese-based creole language in the world'' (2025), which covers local revitalization efforts: \url{https://globalvoices.org/2025/05/08/taiwans-yilan-creole-revitalizing-the-only-japanese-based-creole-language-in-the-world/}
\item TaiwanPlus, ``Meet the Speakers of the World's Last Japanese-Based Creole Language'' (documentary feature): \url{https://www.taiwanplus.com/news/highlights/connected/250402014/meet-the-speakers-of-the-worlds-last-japanese-based-creole-language}
\end{itemize}

\section{Curriculum Development Context}
Yilan Creole has few speakers, mostly elderly, and no recognised place in schools. Existing material is academic and journalistic. Any curriculum would need community-led documentation, and the sources point to Mandarin as the main pressure on the language.
\end{document}